\section{Upper bound on the slice rank}\label{sec:upper}

We now bound the slice rank of the constraint tensor. Define
$T:(\F_3^n)^3\to\F_3$ by $T(x,y,z):=\one[x+y+z=0]$, the indicator that the three
points sum to zero.

\begin{lemma}[Upper bound]\label{lem:upper}
The tensor $T(x,y,z)=\one[x+y+z=0]$ on $(\F_3^n)^3$ admits the polynomial
representation
\begin{align}\label{eq:T-poly}
T(x,y,z)\;=\;\prod_{i=1}^{n}\bigl(1-(x_i+y_i+z_i)^2\bigr),
\end{align}
a reduced polynomial of total degree $2n$, and its slice rank satisfies
$\sr(T)\le 3M_n$, where $M_n$ is as in \Cref{def:monomial-space}.
\end{lemma}

\begin{proof}
The proof has two parts: the polynomial identity Eq.~\eqref{eq:T-poly}, then the slice
decomposition obtained by grouping monomials according to their lowest-degree block
of variables.

\smallskip\noindent
\textbf{Step 1: the polynomial identity.}
For $u\in\F_3$, $\ u^2=1$ if $u\ne 0$ and $u^2=0$ if $u=0$, so $1-u^2=\one[u=0]$.
Applying this coordinatewise with $u=x_i+y_i+z_i$,
\begin{align*}
\prod_{i=1}^{n}\bigl(1-(x_i+y_i+z_i)^2\bigr)
\;\overset{(a)}{=}\;\prod_{i=1}^{n}\one[x_i+y_i+z_i=0]
\;\overset{(b)}{=}\;\one[x+y+z=0]
\;\overset{(c)}{=}\;T(x,y,z),
\end{align*}
where (a) is $1-u^2=\one[u=0]$ per coordinate, (b) is that a product of indicators is
the indicator of the intersection, and (c) is the definition of $T$. Each of the $n$
factors has degree $2$, so the product is reduced of total degree $2n$.

\smallskip\noindent
\textbf{Step 2: monomial grouping.}
Expand Eq.~\eqref{eq:T-poly} into reduced monomials,
\begin{align}\label{eq:T-expand}
T(x,y,z)\;=\;\sum_{(a,b,c)}\gamma_{a,b,c}\,x^a y^b z^c,
\qquad a,b,c\in\{0,1,2\}^n,\ \ \deeg(x^a)+\deeg(y^b)+\deeg(z^c)\le 2n,
\end{align}
the total-degree bound being inherited from Step~1. Since the three block-degrees of
each monomial sum to $\le 2n$, their minimum is $\le 2n/3$:
\begin{align}\label{eq:min-degree}
\min\set{\deeg(x^a),\,\deeg(y^b),\,\deeg(z^c)}\;\le\;\frac{2n}{3}.
\end{align}
Assign each monomial to the block attaining a degree $\le 2n/3$ (ties broken by the
order $x\prec y\prec z$) and collect class by class:
\begin{align}\label{eq:three-groups}
T(x,y,z)
&\;=\;\underbrace{\sum_{\deeg(x^a)\le 2n/3}x^a\,p_a(y,z)}_{=:~T_1}
\;+\;\underbrace{\sum_{\deeg(y^b)\le 2n/3}y^b\,q_b(x,z)}_{=:~T_2}
\;+\;\underbrace{\sum_{\deeg(z^c)\le 2n/3}z^c\,r_c(x,y)}_{=:~T_3},
\end{align}
where $p_a,q_b,r_c$ collect the remaining factors; by Eq.~\eqref{eq:min-degree} each
monomial lands in exactly one class, so Eq.~\eqref{eq:three-groups} is an identity.
Each summand $x^a\,p_a(y,z)$ is a slice in the first coordinate (Eq.~\eqref{eq:slice}
with $f=x^a$), and the monomials $x^a$ with $\deeg(x^a)\le 2n/3$ number exactly
$M_n=\dim V_{\le 2n/3}$ (\Cref{def:monomial-space}); thus $\sr(T_1)\le M_n$, and
likewise $\sr(T_2),\sr(T_3)\le M_n$. By subadditivity \Cref{lem:slice-rank-def:sub},
\begin{align*}
\sr(T)\;\le\;\sr(T_1)+\sr(T_2)+\sr(T_3)\;\le\;3M_n,
\end{align*}
which is the claim. This completes the proof.
\end{proof}

\begin{remark}[Where the threshold $2n/3$ comes from]\label{rem:threshold}
The exponent $2n/3$ in \Cref{def:monomial-space} is exactly $d/3$ for the degree
$d=2n$ of $T$: with three variable blocks summing to degree $\le d$, the minimum
block-degree is $\le d/3$, which is the bound Eq.~\eqref{eq:min-degree}. This is the
slice-rank ($3$-tensor) analogue of the Croot–Lev–Pach matrix bound
$\rank\le 2M_{d/2}$ \cite{crootlevpach2017}, where two variable blocks force a
minimum of $\le d/2$.
\end{remark}
